
\documentclass[11pt]{article}
\usepackage[margin=1in]{geometry}
\usepackage{hyperref}
\usepackage{setspace}
\usepackage[backend=biber, citestyle=authoryear-icomp, bibstyle=authoryear]{biblatex}
\addbibresource{bib.bib}

\onehalfspacing
\setlength{\parindent}{0.5in}

\title{Gendered Scrutiny under Scarcity: First Impressions, Endogenous Information Search, and a Novel Economic Conjoint Framework for Electoral Judgment}
\author{}
\date{\today}

\begin{document}

\maketitle

\section*{Introduction}

Research in political psychology and political behavior has long shown that voters rely on low-information cues when evaluating candidates, especially facial appearance, perceived competence, attractiveness, trustworthiness, and other first-impression signals. At the same time, work in political economy and behavioral decision-making suggests that information use is shaped not only by limited cognition but also by the costs of acquiring and processing additional evidence. Yet these literatures have rarely been integrated in a way that treats information search itself as part of the political judgment process. Most existing studies ask whether appearance affects evaluations, or whether providing new information changes preferences. {\bf Much less is known about whether voters actively choose to remain at the level of first impressions, when they decide to investigate further, and whether those search decisions are themselves structured by gender}.

This paper argues that electoral judgment is not only a matter of which candidate traits voters prefer, but also of how much effort they are willing to expend to learn about different candidates. That question is especially important in the case of women in politics. A large literature suggests that female candidates are often evaluated through gendered expectations and may face distinct standards of competence, likability, and suitability for office. But the existing evidence focuses mainly on vote choice, candidate evaluations, or stereotype activation. {\bf It tells us whether gender matters for outcomes, but not whether gender shapes the process of attention allocation itself}. We contend that one of the key mechanisms of unequal political judgment may lie precisely there: female candidates may elicit greater information demand because respondents feel a stronger need to verify, resolve, or reconcile their initial impressions. In that sense, political judgment is distributed \emph{unevenly} across candidates.

Our core theoretical claim is that first impressions generate different levels of subjective confidence, and that this confidence governs the demand for additional information. When respondents believe that a candidate's face already provides a sufficiently clear signal, they may rationally decline to spend scarce resources on further investigation: low search may reflect efficient economizing if first impressions are informative, but it may also reflect overconfidence if respondents stop too early and rely too heavily on facial heuristics. The gendered implication is especially important: {\bf if female candidates are judged under stricter or less stable priors, respondents may be less willing to trust first impressions alone and more likely to invest in additional search}. The result is a stronger demand for information not because women inherently provide less information, but because they are more heavily scrutinized in the act of political evaluation.

To study this process, we introduce a novel \emph{economic conjoint experiment} implemented in oTree (see prototype \href{https://econ-conjoint-project.onrender.com/}{here}). The design is novel because it combines the multidimensional choice structure of conjoint analysis with an explicitly economic search environment in which information is scarce, costly, and endogenously acquired. Rather than presenting respondents with a fully revealed bundle of candidate attributes, we allow them to uncover information selectively by spending points from a limited budget. This transforms the conjoint from a passive stated-preference task into a sequential decision problem in which \emph{search behavior} itself becomes theoretically meaningful. Respondents first encounter real candidate photographs in a low-information setting, form an initial judgment, and then decide whether to pay to reveal additional profile attributes before updating their choice. The experiment therefore identifies not only which candidate features affect electoral judgments, but also which features respondents consider worth learning about, under what conditions, and for whom.

This methodological contribution is substantial. Standard conjoint experiments are powerful tools for estimating the causal effects of profile attributes on selection probabilities, but they usually treat respondents as passive integrators of information that the researcher has already chosen to display. As a result, they reveal little about how subjects allocate attention, whether they satisfice, or when they decide that they have seen enough. Our economic conjoint experiment addresses that limitation directly. By imposing a binding search budget and recording click-level information acquisition, it produces revealed-preference measures of attention allocation, confidence, and deliberation. 

We expect respondents to spend fewer points when faces generate strong confidence signals and more when visual cues are ambiguous or incongruent. But we also expect this logic to be gendered: female candidates should trigger greater information demand, consistent with a politics of asymmetric scrutiny in which women are evaluated under a higher burden of verification. By combining real candidate images, endogenous information acquisition, repeated decision tasks, and budget-constrained search, the paper offers a unified framework for studying how first impressions, gender, and costly learning jointly structure electoral judgment.

%Link to the prototype \href{https://econ-conjoint-project.onrender.com/}{here}.

\printbibliography

\end{document}

\section*{Methodological Advantages} 

The experimental design offers several behavioral and methodological advantages for studying electoral judgment under conditions of scarcity. 

Respondents first encounter candidates in a low-information environment defined only by facial appearance, and must then decide whether to spend points to reveal additional attributes---age, occupation, and party---before making their judgment. This feature turns information acquisition into an observable behavioral outcome rather than treating it as an unobserved assumption. 

The design therefore makes it possible to distinguish reliance on first impressions from active verification, and to identify when respondents consider the initial visual signal sufficient and when they judge that further investigation is worth the cost. Because the additional information is standardized across candidates, differences in search behavior can be interpreted more cleanly as differences in confidence, scrutiny, or perceived ambiguity rather than as artifacts of unequal text length, rhetorical style, or linguistic complexity. This structure also improves on conventional conjoint experiments by opening the black box between candidate exposure and final choice. 

Standard conjoint designs usually reveal all attributes at once and therefore identify how respondents react to bundles of information that have already been selected and displayed by the researcher. By contrast, the present design records whether respondents choose to uncover age, occupation, and party, how much they are willing to pay to do so, and how those decisions shape their final judgment. In this sense, the experiment generates revealed-preference measures of attention allocation under scarcity. It thus permits a more direct test of whether low-information judgments reflect efficient economizing, false confidence, or asymmetric standards of evaluation. The repeated-task structure further strengthens the design because it allows the researcher to observe whether respondents develop stable search strategies over time and whether costly information acquisition varies systematically across candidates. More substantively, age, occupation, and party are especially useful attributes because they speak directly to several of the central heuristics voters use in candidate evaluation. Occupation may signal competence, prestige, and social status; party provides an immediate political anchor; and age may shape perceptions of experience and suitability for office. 

Allowing respondents to reveal these cues selectively makes it possible to study not only which attributes matter, but also for whom they are judged necessary. The design is therefore well suited to testing whether some candidates are evaluated under a higher burden of verification than others, and whether electoral judgment is shaped as much by unequal scrutiny and differential confidence as by preferences over candidate traits themselves.


